\documentclass[11pt]{article}
\usepackage[margin=1in]{geometry}
\usepackage{amsmath,amssymb,amsthm}
\usepackage{xurl}
\usepackage{hyperref}
\sloppy
\let\oldus\_ \renewcommand{\_}{\oldus\allowbreak}
\newtheorem{theorem}{Theorem}
\newtheorem{lemma}[theorem]{Lemma}
\theoremstyle{definition}
\newtheorem{definition}[theorem]{Definition}
\theoremstyle{remark}
\newtheorem{remark}[theorem]{Remark}

\title{Disproof of conjecture 00000000086:\\ $\log|B(m,p)|\not\sim\frac{(m-1)^2}{2}\log p$, already for $m=2$}
\author{Jackmeson1}
\date{2026-10-05}

\begin{document}
\maketitle

\section{The conjecture and the reading}

\textbf{Statement (English, verbatim).} Definition: The free Burnside group $B(m,p)$ denotes the
largest $m$-generated group satisfying $x^p = 1$. Conjecture: $\log|B(m,p)|$ is asymptotic to
$(m-1)^2/2\cdot\log p$; and this asymptotic holds as $p$ tends to infinity. (The Chinese text
states the same.)

\textbf{Reading.} For fixed $m$, as $p\to\infty$,
\[
  \log|B(m,p)| \;\sim\; \tfrac{(m-1)^2}{2}\,\log p ,
\]
where $u\sim v$ means $u-v=o(v)$ (Mathlib \texttt{Asymptotics.IsEquivalent}). The rank $m$ is a
free parameter, so the conjecture asserts this for every $m$; we take the weakest such form, every
$m\ge 2$, and refute it at $m=2$. That refutes the universal claim for any range of $m$ containing
$2$. We also refute the individual claims for $m=1$ and $m=3$.

\textbf{Conventions.} $\log$ is the natural logarithm. The base does not matter: changing it
multiplies both sides by the same positive constant, and $\sim$ is invariant under this. The
exponent $p$ ranges over the natural numbers. ``$p\to\infty$'' is the filter \texttt{atTop}. We
also treat $p\to\infty$ through the primes, and in fact any filter $l\le\texttt{atTop}$ with
$l\ne\bot$. $|G|$ is \texttt{Nat.card G}, which is $0$ when $G$ is infinite (Mathlib's junk value),
so $\log|G|=\log 0=0$ (Mathlib's \texttt{Real.log 0 = 0}). We refute two readings of $\log|B|$ for
possibly infinite $B$:
\begin{itemize}
\item \emph{Strict reading} (\texttt{Claim}): $B(m,p)$ is finite for all large $p$ (so that
  $\log|B(m,p)|$ is a real number), and the asymptotic holds.
\item \emph{Junk reading}: the asymptotic holds for $\log(\texttt{Nat.card})$, with no finiteness
  requirement.
\end{itemize}
A reading in which $\log|B|=+\infty$ for infinite $B$ fails at every such $p$ in the same way as the
strict reading. We do not refute the claim for an individual $m\ge 4$ (see Remark~\ref{rem:m4}).

\section{The free Burnside group (as formalized)}

\begin{definition}[\texttt{B}, \texttt{gen}]
$B(m,p)=F_m/N$, where $F_m$ is the free group on $\{0,\dots,m-1\}$ (\texttt{FreeGroup (Fin m)}) and
$N$ is the normal closure of the set $\{g^p : g\in F_m\}$ (\texttt{powers m p}). The generators are
the images $x_i$ of the free generators.
\end{definition}

\begin{definition}[\texttt{IsLargest m p G}]
A group $G$ is a \emph{largest $m$-generated group satisfying $x^p=1$} if (i) $G$ is generated by
some $m$ elements, (ii) $x^p=1$ for every $x\in G$, and (iii) for every group $H$ that satisfies (i)
and (ii), there is a surjective homomorphism $G\to H$.
\end{definition}

Wikipedia's article on the Burnside problem defines $B(m,n)$ by this universal property: given
any group $G$ with $m$ generators and of exponent $n$, ``there is a unique homomorphism from
$B(m, n)$ to $G$ that maps the ith generator xi of $B(m, n)$ into the ith generator gi of $G$''.
We prove exactly this property for our $B(m,p)$:

\begin{lemma}\label{lem:univ}
\begin{enumerate}
\item $x^p=1$ for every $x\in B(m,p)$ (\texttt{B\_pow\_eq\_one}).
\item $B(m,p)$ is generated by $x_0,\dots,x_{m-1}$ (\texttt{closure\_range\_gen}).
\item For every group $H$ with $h^p=1$ for all $h\in H$, and every map $f$ from
  $\{0,\dots,m-1\}$ to $H$, there is a homomorphism $\mathrm{lift}\,f:B(m,p)\to H$ with
  $x_i\mapsto f(i)$ (\texttt{lift}, \texttt{lift\_gen}). It is unique: two homomorphisms that agree
  on the $x_i$ are equal (\texttt{hom\_ext}). If the $f(i)$ generate $H$, then $\mathrm{lift}\,f$ is
  surjective (\texttt{lift\_surjective}).
\item Hence $B(m,p)$ satisfies \texttt{IsLargest m p} (\texttt{B\_isLargest}).
\end{enumerate}
\end{lemma}

\begin{proof}
(1) Every $x$ is the image of some $g\in F_m$, and $g^p\in N$. (2) The free generators generate
$F_m$, and the quotient map is surjective. (3) $N$ lies in the kernel of the free-group lift of
$f$, because $f$-images of $p$-th powers are $p$-th powers in $H$, hence trivial. Uniqueness holds
because a homomorphism out of $F_m/N$ is determined on $F_m$, and on $F_m$ by the free generators.
Surjectivity: the range is a subgroup that contains every $f(i)$.
\end{proof}

All results below are proved for an arbitrary family of groups satisfying \texttt{IsLargest}, so
they do not depend on our particular construction. Lemma~\ref{lem:univ} shows that this hypothesis
is satisfied.

\section{The proof}

\begin{lemma}[\texttt{E\_pow\_eq\_one}, \texttt{closure\_range\_egen}, \texttt{card\_E}]
Let $E(m,p)=(\mathbb{Z}/p)^m$, written multiplicatively (\texttt{Multiplicative (Fin m -> ZMod p)}).
Then $x^p=1$ for every $x$. For $p\ge1$, $E(m,p)$ is generated by the $m$ standard basis vectors.
Also $|E(m,p)|=p^m$.
\end{lemma}

\begin{proof}
$p\cdot a=0$ in $\mathbb{Z}/p$. For the second claim, $x=\sum_i x_i e_i$, a product of powers of
the basis vectors.
\end{proof}

\begin{lemma}[\texttt{card\_ge}, \texttt{finite\_of\_isLargest\_one}]\label{lem:card}
Let $G$ satisfy \texttt{IsLargest m p}, with $p\ge1$. If $G$ is finite, then $|G|\ge p^m$. If
$m=1$, then $G$ is finite.
\end{lemma}

\begin{proof}
By (iii), $G$ maps onto $E(m,p)$, so $|G|\ge|E(m,p)|=p^m$. If $m=1$, then $G$ is generated by one
element $g$ with $g^p=1$, so $G=\langle g\rangle$ is finite.
\end{proof}

\begin{lemma}[\texttt{not\_isEquivalent\_of\_gap}]
Let $l$ be a filter on $\mathbb{N}$ with $l\ne\bot$, and let $\eta>0$. If
$\eta|v(p)|<|u(p)-v(p)|$ for $l$-eventually all $p$, then $u\not\sim v$ along $l$.
\end{lemma}

\begin{proof}
$u\sim v$ would give $|u(p)-v(p)|\le\eta|v(p)|$ eventually. Both eventualities hold at some $p$,
which is a contradiction.
\end{proof}

\begin{lemma}[\texttt{gap}]\label{lem:gap}
Let $m\in\{1,2,3\}$, $p\ge 2$, $G$ satisfy \texttt{IsLargest m p}, $L=\log p>0$,
$\kappa=(m-1)^2/2\in\{0,\frac12,2\}$ and $v=\kappa L$. Then
$\tfrac14|v|<|\log|G|-v|$ (with $\log|G|=0$ if $G$ is infinite).
\end{lemma}

\begin{proof}
If $G$ is finite, Lemma~\ref{lem:card} gives $\log|G|\ge mL$, so
$\log|G|-v\ge(m-\kappa)L$. Here $m-\kappa\in\{1,\frac32,1\}$, which is $>\frac14\kappa\in\{0,\frac18,\frac12\}$.
If $G$ is infinite, then $m\in\{2,3\}$ by Lemma~\ref{lem:card}, $\log|G|=0$, and
$|0-v|=\kappa L>\frac14\kappa L$ since $\kappa>0$.
\end{proof}

\begin{theorem}[\texttt{not\_asymptotic\_of\_isLargest}, \texttt{burnside\_not\_asymptotic},
\texttt{burnside\_not\_claim}]
Let $l\le\texttt{atTop}$ be a filter on $\mathbb{N}$ with $l\ne\bot$, and let $m\in\{1,2,3\}$.
For every family $(G_p)$ of groups with \texttt{IsLargest m p} $G_p$ for all $p$, and in particular
for $G_p=B(m,p)$, the function $\log|G_p|$ (junk value $0$ for infinite $G_p$) is not asymptotic to
$\frac{(m-1)^2}{2}\log p$ along $l$. Consequently the strict reading \texttt{Claim l m (B m)} is
false.
\end{theorem}

\begin{proof}
$p\ge2$ holds $l$-eventually, since $l\le\texttt{atTop}$. Apply the gap lemma with $\eta=\frac14$ and
Lemma~\ref{lem:gap}. The strict reading contains the asymptotic as a conjunct.
\end{proof}

\begin{theorem}[\texttt{conjecture86\_false}, \texttt{conjecture86\_false\_atTop},
\texttt{conjecture86\_false\_primes}]
\texttt{Conjecture l}, defined as ``for every $m\ge2$: $B(m,p)$ is eventually finite along $l$ and
$\log|B(m,p)|\sim\frac{(m-1)^2}{2}\log p$ along $l$'', is false for every filter $l\le\texttt{atTop}$
with $l\ne\bot$. This includes $l=\texttt{atTop}$ (all integers $p\to\infty$) and
$l=\texttt{atTop}\sqcap\mathcal{P}\{p\text{ prime}\}$ (primes $p\to\infty$). The filter of primes is
not $\bot$, because there are infinitely many primes (\texttt{primes\_neBot}).
\end{theorem}

\begin{proof}
Take $m=2$. Concretely, if $B(2,p)$ is finite then $|B(2,p)|\ge p^2$, so
$\log|B(2,p)|/(\frac12\log p)\ge4$ for every $p\ge2$. The ratio cannot tend to $1$.
\end{proof}

\begin{remark}\label{rem:m4}
For $m\ge4$ we have $m\le(m-1)^2/2$, and the abelian quotient no longer separates the two sides.
We do not formalize the $m\ge4$ instances. As context only (not used in the proof): the Wikipedia
article on the Burnside problem states that Ivanov's theorem, together with the Novikov--Adian
theorem, implies infiniteness for all $m>1$ and $n\ge 2^{48}$ (the exponent is a superscript in the
article; plain-text extraction renders it as ``248''). Under the strict reading, this
would refute every instance $m\ge2$.
\end{remark}

\section{Lean formalization and axiom audit}

The file \texttt{lean/Conjecture86/Basic.lean} (235 lines, Mathlib v4.33.1, only
\texttt{import Mathlib}) contains all the definitions and theorems named above. The main
statement is:
\begin{verbatim}
theorem conjecture86_false (l : Filter Nat) [l.NeBot]
    (hl : l <= Filter.atTop) : Not (Conjecture l)
def Conjecture (l : Filter Nat) : Prop :=
  forall m : Nat, 2 <= m -> Claim l m (B m)
def Claim (l : Filter Nat) (m : Nat) (G : Nat -> Type) : Prop :=
  Filter.Eventually (fun p => Finite (G p)) l /\
    Asymptotics.IsEquivalent l (logOrder G) (target m)
\end{verbatim}
This is an ASCII rendering of the Lean source, which uses Unicode notation. Here
\texttt{target m p = ((m : Real) - 1)\^{}2 / 2 * Real.log p} and
\texttt{logOrder G p = Real.log (Nat.card (G p))}. \texttt{lake build} succeeds with no
\texttt{sorry} and no \texttt{native\_decide}. \texttt{\#print axioms} gives
\texttt{[propext, Classical.choice, Quot.sound]} for each of \texttt{conjecture86\_false},
\texttt{conjecture86\_false\_atTop}, \texttt{conjecture86\_false\_primes},
\texttt{burnside\_not\_asymptotic}, \texttt{not\_asymptotic\_of\_isLargest},
\texttt{B\_isLargest} and \texttt{hom\_ext}.

\section*{Source}
Wikipedia, ``Burnside problem'', \url{https://en.wikipedia.org/wiki/Burnside_problem} (retrieved
2026-10-05). It is used for the definition of $B(m,n)$ and for the context in
Remark~\ref{rem:m4}.

\end{document}
